\documentclass[14pt, a4paper]{extreport}

\usepackage{kpi-lab}
% Defined by subject
\newcommand{\CourseTitle}{Бізнес-аналіз в ІТ}
\newcommand{\Teacher}{ст. в. Міщенко Л. Д.}
\newcommand{\ReportYear}{2026}
% Defined by Student
\newcommand{\StudentName}{Ковальов Олександр}
\newcommand{\StudentGroup}{ІМ-51мн}
\newcommand{\Variant}{4}
\newcommand{\CourseNumber}{2}
% Defined by Submission
\newcommand{\Topic}{Збір вимог (інтерв’ю, опитування)}
\newcommand{\LabNumber}{2}

\begin{document}
	\pagenumbering{gobble}
	
	\begin{titlepage}
		\begin{center}
			\normalsize
			Національний технічний університет України\\
			«Київський політехнічний інститут імені Ігоря Сікорського» \\[0.5em]
			Факультет інформатики та обчислювальної техніки\\
			Кафедра обчислювальної техніки
			
			\vfill
			
			{\large \textbf{ЗВІТ}\\[0.5em]}
			{з лабораторної роботи №\LabNumber \\}
			{з дисципліни <<\CourseTitle>>}
			
			\vspace{2em}
			
			{\textbf{<<\Topic>>}}
			
			\vspace{1em}
			
			{\textbf{Варіант №\Variant}}
			
			\vfill
			
			\begin{flushright}
				Виконав: \\ студент \CourseNumber{} курсу, групи \StudentGroup \\
				\StudentName \\[1em]
				Перевірила: \\ \Teacher
			\end{flushright}
			
			\vfill
			
			\begin{flushright}
				Дата здачі: \DTMtoday
			\end{flushright}
			
			\vspace{5em}
			КИЇВ -- \ReportYear
		\end{center}
	\end{titlepage}
	
	\clearpage
	\pagenumbering{arabic}
	\setcounter{page}{2}
	
	\subsection*{Завдання}
	
	\begin{enumerate}
		\item \textbf{Інтерв'ю}. Проведіть інтерв'ю з уявним стейкхолдером: 10-12 питань, документування відповідей та висновки аналітика. Питання повинні бути глибинні, які б давали змістовні відповіді.
		
		\item \textbf{Опитування}. Створіть опитування з 10-12 запитань різного типу через Google Forms або \underline{текстовий документ}. Треба для поверхневих та загальних запитань. Оцінюється адекватність питань, баланс відкритих і закритих форматів.
		
		\item \textbf{User Stories}. Напишіть 8-12 історій користувачів за стандартним форматом на основі зібраних даних. Оцінюється повнота, коректний формат, чіткий зв'язок із зібраними даними.
		
		\item \textbf{Use-Case Diagram}. Створіть діаграму використання з акторами, use cases та межами системи. Повинні бути чіткі актори, логічні сценарії, правильна структура діаграми.
	\end{enumerate}
	
	Ключовими висновками повинні бути:
	
	\begin{enumerate}
		\item \textbf{Мета} -- формалізувати вимоги стейкхолдерів через структуровані методи збору інформації.
		
		\item \textbf{Результат} -- завдання, описані вище, які становлять основу для подальшого проектування системи.
		
		\item \textbf{Наступний крок} -- використання артефактів для розробки архітектури та детального дизайну рішення.
	\end{enumerate}
	
	\section*{Хід роботи}
	
	\textbf{1. Інтерв'ю зі стейкхолдером}
	
	Для глибинного розуміння потреб було проведено інтерв'ю з уявним стейкхолдером -- власником приватного будинку, який активно використовує зарядні станції та планує встановлення сонячних панелей. 
	
	\begin{enumerate}
		\item \textbf{<<Як зараз організовано процес резервного живлення у вашому будинку?>>} Респондент відповів, що використовує зарядну станцію, яку доводиться підключати до щитка вручну при кожному відключенні. 
		
		\item \textbf{<<Які найбільші незручності виникають під час знеструмлень?>>}. Було зазначено, що головною проблемою є необхідність бути присутнім вдома для перемикання рубильників та постійний стрес через неточні графіки.
		
		\item \textbf{<<Скільки часу ви витрачаєте на моніторинг ситуації з електроенергією?>>}. Стейкхолдер зізнався, що перевіряє Telegram-канали обленерго по кілька разів на день.
		
		\item \textbf{<<Чи стикалися ви з фінансовими втратами через відключення?>>}. Він згадав про зіпсовані продукти в холодильнику через розряджену вночі станцію.
		
		\item \textbf{<<Яким ви бачите ідеальне рішення для керування живленням?>>}. Респондент описав систему, яка б сама перемикала джерела і надсилала сповіщення на телефон.
		
		\item \textbf{<<Наскільки важливим є віддалене керування через смартфон?>>}. Була отримана відповідь, що це абсолютно критично, оскільки часто доводиться залишати будинок на тривалий час.
		
		\item \textbf{<<Які сценарії автоматизації для вас найбільш актуальні?>>}. Стейкхолдер виділив заряджання акумуляторів виключно вночі за низьким тарифом.
		
		\item \textbf{<<Як ви ставитеся до прогнозування на основі погоди?>>}. Він зазначив, що це було б дуже корисно для майбутніх сонячних панелей.
		
		\item \textbf{<<Який бюджет ви вважаєте прийнятним для такої автоматизації?>>}. Стейкхолдер оцінив готовність витратити близько 500 доларів за надійний розумний контролер.
		
		\item \textbf{<<Чи є побоювання щодо безпеки системи?>>}. Респондент висловив занепокоєння щодо можливих замикань при автоматичному перемиканні реле.
	\end{enumerate}

	\textbf{Висновок аналітика:} Інтерв'ю підтверджує гіпотезу, що користувачі втомилися від ручного керування. Ключовою цінністю продукту XaiVolt має стати повна автономність перемикань, прозора аналітика в застосунку та безкомпромісна безпека обладнання.
	
	\textbf{2. Опитування користувачів}
	
	Для валідації проблеми на ширшій аудиторії було розроблено опитування з десяти запитань різного типу. Перше питання з'ясовувало тип житла респондента (закритий тип: квартира, будинок, комерція). Друге питання стосувалося наявного обладнання (множинний вибір: генератор, зарядна станція, інвертор). Третє питання пропонувало оцінити рівень дискомфорту від ручного перемикання за шкалою від 1 до 10. Четверте запитувало про середню кількість відключень на тиждень. П'яте питання стосувалося того, чи бувають респонденти відсутні вдома під час знеструмлень (так/ні). Шосте питання пропонувало обрати найбажанішу функцію системи (множинний вибір: віддалене керування, статистика, автоперемикання). Сьоме запитання було відкритим і просило описати найбільший страх при використанні резервного живлення. Восьме питання з'ясовувало готовність довірити системі керування потужними приладами, такими як бойлер (шкала згоди). Дев'яте питання стосувалося бажаного формату сповіщень (закритий тип: push, SMS, Telegram). Десяте питання було відкритим і збирало пропозиції щодо ідеального мобільного застосунку. Такий баланс дозволив би зібрати як кількісні метрики щодо обладнання, так і якісні інсайти щодо побоювань.
	
	\textbf{3. User Stories}
	
	На основі зібраних даних було сформовано набір історій користувачів. 
	
	\begin{enumerate}
		\item Як власник будинку, я хочу налаштувати автоматичне перемикання на резервне живлення, щоб не робити це вручну в темряві.
		
		\item Як користувач системи, я хочу отримувати push-сповіщення про перехід на роботу від акумулятора, щоб бути в курсі стану мережі.
		
		\item Як економний споживач, я хочу, щоб система заряджала станцію вночі за пільговим тарифом, щоб зменшити витрати.
		
		\item Як власник сонячних панелей, я хочу бачити прогноз генерації на основі погоди, щоб планувати використання енергоємних приладів.
		
		\item Як працююча людина, я хочу віддалено переглядати рівень заряду батарей через застосунок, щоб знати, чи вистачить енергії до мого повернення.
		
		\item Як користувач, я хочу мати можливість вручну вимкнути живлення певних розеток через смартфон, щоб зберегти заряд.
		
		\item Як власник бізнесу, я хочу бачити детальну статистику споживання за місяць, щоб оптимізувати витрати.
		
		\item Як системний адміністратор, я хочу отримувати сповіщення про перегрів інвертора, щоб вчасно запобігти аварії.
		
		\item Як користувач екосистеми, я хочу інтегрувати систему з графіками відключень ДТЕК, щоб контролер міг превентивно дозарядити батареї.
		
		\item Як обережний користувач, я хочу мати фізичну кнопку аварійного відключення, щоб швидко знеструмити будинок у разі небезпеки.
	\end{enumerate}

	\textbf{4. Use-Case Diagram}
	
	Для візуалізації функціональних вимог розроблено Use-Case діаграму. Основними акторами виступають <<Користувач (Власник)>>, який взаємодіє із системою через мобільний застосунок, <<Апаратний контролер>>, що виконує фізичні перемикання, <<Зовнішні API>>, які надають дані про погоду та графіки обленерго, <<Електрик-інсталятор>>, відповідальний за первинний монтаж, та <<Системний адміністратор>>, який здійснює технічну підтримку і діагностику обладнання. Межі системи охоплюють комплексну програмну логіку та взаємодію з контролером. Модель включає складні сценарії із залежностями. Зокрема, реалізовано відношення обов'язкового включення (\texttt{include}), де управління розетками та налаштування сценаріїв користувачем вимагають попередньої авторизації в системі, а первинне налаштування контролера інсталятором завжди супроводжується калібруванням датчиків. Відношення розширення (\texttt{extend}) дозволяють опціонально доповнювати базовий моніторинг стану мережі переглядом детальної статистики, тоді як стандартне автоматичне перемикання реле може ініціюватися прецедентом аварійного відключення у разі виявлення замикання. Водночас зовнішні сервіси безперервно постачають дані, необхідні для синхронізації графіків відключень та прогнозування сонячної генерації.
	
	\image{14cm}{use-case}
	
	\subsection*{Висновок}
	
	Головною метою цієї роботи було формалізувати вимоги стейкхолдерів через структуровані методи збору інформації. Завдяки проведеному інтерв'ю та підготовленому опитуванню вдалося глибше зрозуміти реальні проблеми користувачів, зокрема їхню втому від ручного керування та непередбачуваності графіків. Результатом виконання завдань стали деталізовані історії користувачів та діаграма прецедентів, які чітко окреслюють межі та функціонал майбутнього продукту. Наступним кроком стане використання створених артефактів для розробки архітектури бази даних та детального технічного дизайну програмно-апаратного комплексу XaiVolt.
	
\end{document}